\section{РАЗРАБОТКА ЛАБОРАТОРНЫХ РАБОТ}

На базе разработанного учебно-лабораторного стенда был сформирован комплект из шести лабораторных работ, охватывающих ключевые аспекты технологии Bluetooth: от базовых процедур установления соединения до анализа трафика и сравнения производительности различных протоколов. Каждая работа включает теоретическую справку, пошаговое руководство по выполнению эксперимента и контрольные вопросы для закрепления материала. Ниже приведено краткое описание содержания и целей каждой лабораторной работы.

\subsection{Лабораторная работа №1. Исследование процессов обнаружения, сопряжения и
установления соединения в сетях Bluetooth}
Данная работа направлена на изучение фундаментальных механизмов взаимодействия устройств в сетях Bluetooth Classic (BR/EDR). Студенты осваивают работу с официальным стеком BlueZ в ОС Linux, используя утилиты \texttt{bluetoothctl} и \texttt{btmon}.
\textbf{Основные изучаемые вопросы:}
\begin{itemize}[nosep]
    \item Процедуры поиска устройств (Inquiry) и вызова (Page);
    \item Механизмы безопасного сопряжения (Secure Simple Pairing, SSP) ~\cite{bt_security_ssp} и связывания (Bonding);
    \item Структура HCI-пакетов, отвечающих за аутентификацию и шифрование канала;
    \item Установление логического соединения на уровне L2CAP.
\end{itemize}
В ходе эксперимента студенты выполняют сканирование эфира, проводят сопряжение с устройством ESP32 и анализируют логи \texttt{btmon} для визуализации этапов handshake-процедуры.

\subsection{Лабораторная работа №2. Исследование архитектуры пикосети Bluetooth и организация обмена данными через RFCOMM}
Работа посвящена изучению топологии «пикосеть» и реализации обмена данными в режиме «мастер–ведомые». Студенты программируют два модуля ESP32 в режиме SPP-серверов и разрабатывают Python-скрипт на ПК, выступающий в роли мастера.
\textbf{Основные изучаемые вопросы:}
\begin{itemize}[nosep]
    \item Принципы синхронизации устройств в пикосети и технология FHSS;
    \item Протокол RFCOMM как эмуляция последовательного порта RS-232 ~\cite{rfcomm_spp_guide};
    \item Мультиплексирование соединений и управление потоком данных;
    \item Программирование Bluetooth-сокетов в Python.
\end{itemize}
Эксперимент включает создание параллельных RFCOMM-соединений с двумя узлами и проверку целостности передаваемых данных.

\subsection{Лабораторная работа №3. Анализ трафика Bluetooth с использованием Wireshark}
Целью работы является получение навыков анализа трафика с помощью Wireshark ~\cite{wireshark_bt_doc}. Студенты учатся перехватывать сырой трафик Bluetooth-адаптера и декодировать пакеты различных уровней стека.
\textbf{Основные изучаемые вопросы:}
\begin{itemize}[nosep]
    \item Структура пакетов Bluetooth (Access Code, Header, Payload);
    \item Фильтрация трафика по типам событий HCI, протоколам L2CAP, RFCOMM и SDP;
    \item Анализ процедур Service Discovery (обнаружение сервисов);
    \item Визуализация аудио-трафика (A2DP) и данных профилей HID/SPP.
\end{itemize}
Студенты захватывают трафик процедур Inquiry, Pairing и передачи файлов, затем составляют отчет с расшифровкой ключевых полей пакетов.

\subsection{Лабораторная работа №4. Организация удаленного выполнения команд через Bluetooth-мост на базе ESP32}
В данной работе реализуется архитектура «Клиент–Шлюз–Сервер», где модуль ESP32 выступает в роли прозрачного моста (Bridge) между двумя персональными компьютерами.
\textbf{Основные изучаемые вопросы:}
\begin{itemize}[nosep]
    \item Одновременная работа ESP32 в режимах SPP-клиента и SPP-сервера;
    \item Проблемы буферизации и управления потоком (Flow Control) при ретрансляции данных;
    \item Механизмы удаленного выполнения команд (Remote Code Execution) через subprocess в Python;
    \item Оценка задержек (Latency) при прохождении данных через промежуточный узел.
\end{itemize}
Студенты настраивают цепочку передачи команд от ПК-Клиента к ПК-Серверу через ESP32 и измеряют время отклика системы.

\subsection{Лабораторная работа №5. Исследование архитектуры BLE и реализация кастомных профилей GATT}
Работа фокусируется на технологии Bluetooth Low Energy (BLE). Студенты разрабатывают прошивки для двух модулей ESP32, реализующих роли Central (клиент) и Peripheral (сервер), и создают пользовательский GATT-сервис.
\textbf{Основные изучаемые вопросы:}
\begin{itemize}[nosep]
    \item Отличия стека BLE от Classic Bluetooth (Advertising, Connection Interval);
    \item Иерархия данных GATT: Профиль $\rightarrow$ Сервис $\rightarrow$ Характеристика $\rightarrow$ Дескриптор;
    \item Использование UUID для идентификации сервисов;
    \item Механизмы уведомлений (Notify) и индикаций (Indicate).
\end{itemize}
В ходе эксперимента создается кастомный сервис с характеристиками READ/WRITE/NOTIFY и организуется двунаправленный обмен данными между двумя микроконтроллерами.

\subsection{Лабораторная работа №6. Сравнение скоростных характеристик протоколов Bluetooth (BR, EDR, BLE)}
Заключительная работа носит исследовательский характер и посвящена измерению реальной пропускной способности (Throughput) различных режимов работы Bluetooth.
\textbf{Основные изучаемые вопросы:}
\begin{itemize}[nosep]
    \item Влияние схем модуляции (GFSK, $\pi/4$-DQPSK, 8DPSK) на скорость передачи;
    \item Зависимость пропускной способности от размера пакета (MTU/Payload);
    \item Накладные расходы протоколов верхних уровней (RFCOMM и ATT);
    \item Сравнение эффективности BR/EDR и BLE при передаче больших объемов данных.
\end{itemize}
Студенты проводят серию замеров скорости передачи данных для разных размеров блоков (CHUNK) в Classic Bluetooth и разных значений MTU в BLE, строят сравнительные графики и делают выводы о применимости технологий для различных сценариев (IoT и Multimedia).

Таким образом, разработанный комплект лабораторных работ позволяет студентам последовательно освоить весь спектр технологий Bluetooth, от физического уровня до прикладных профилей, используя созданный лабораторный стенд.
